\section{Research Positions}

\setlength{\arrayrulewidth}{.01pt}
\arrayrulecolor[HTML]{C1C1C1}
\setlength{\tabcolsep}{8pt}
\setlength{\LTpre}{0pt}

%\begin{tabular}{r|p{15.75cm}}
\begin{longtable}{r|p{15.6cm}}
    % 2024 \textendash
    % & \textbf{Microbiology Research Assistant} \\ 
    % \emph{Present}
    % &\footnotesize{\emph{Center for Biofilm Engineering}}
    % \\&\\
    
    \small{\textsc{2022}} & \textbf{North Carolina Center Lead and Communications Fellow} \\ 
    \small{\textendash ~\textsc{2024}}
    &\footnotesize{\emph{NOAA National Centers for Environmental Information, NASA DEVELOP}}\\
    &\footnotesize{Designed and managed six research projects applying NASA Earth observations to environmental management decisions in collaboration with 14 academic, government, and non-profit organizations. Supervised six project teams of 3–5 researchers each. Supported Software Carpentry Workshops on Python and Git. Managed operations of NASA DEVELOP's NOAA NCEI Office and supported career-building for early-career scientists.}
    % Communicate project results to local communities in collaboration with academic, government, and non-profit partners.}
    % Supported Software Carpentry Workshops on geospatial analysis using Python and version control with Git. Managed day-to-day operations of NASA DEVELOP's NOAA NCEI Office, and organized career-building and networking events for early-career scientists.
    \\&\\

    \small{\textsc{2021}} & \textbf{Ecological Forecasting Researcher} \\
    \small{\textendash ~\textsc{2022}}
    &\footnotesize{\emph{NASA Ames Research Center \& University of Georgia Center for Geospatial Research, NASA DEVELOP}}\\
    &\footnotesize{Collaborated with federal, academic, and non-profit agencies to model present and future Venus flytrap habitat and monitor spongy moth defoliation using remote sensing imagery and environmental data. }
    \\&\\

    %\textsc{2019} 
    %&\textbf{Undergraduate Mycology Researcher} \\
    %\textendash ~\textsc{2020}
    %&\footnotesize{\emph{Depts. of Biology, University of North Carolina Asheville and Warren Wilson College}}\\
    %&\footnotesize{Investigated \textit{Tsuga canadensis} decline due to \textit{Adelges tsugae} infestation by comparing seedling growth and ectomycorrhizal colonization rates and taxonomic richness between healthy and declining hemlock stands.}
    %Results published in \textit{American Midland Naturalist}.}
    %\\&\\

    %\textsc{2019} & \textbf{Phenology Research Assistant} \\
    %\textendash ~\textsc{2020}
    %&\footnotesize{\emph{Dept. of Biology, University of North Carolina Asheville}}\\
    %&\footnotesize{Tracked the effects of climate variability on Southern Appalachian perennial plant phenology and monitored regrowth of plant communities following invasive plant removal
%}
    %\\&\\

    \small{\textsc{2019}} 
    &\textbf{Research Fellow} \\
    \small{\textendash ~\textsc{2020}}
    &\footnotesize{\emph{McCullough Institute}}\\
    &\footnotesize{Examined \textit{Spiraea virginiana} reproduction and physiology to assess barriers to species' survival comparing reproductive effort and output among populations and testing pollen and seed viability through selfing, inbreeding, and outbreeding treatments. Facilitated plant education workshops for pre-K\textendash middle school students. Collected and propagated \textit{S. virginiana} for reintroduction and reintroduced only known member of extirpated \textit{S. virginiana} population maintained in Arnold Arboretum Living Collections to Botanical Gardens at Asheville.
}
\\&\\
    \small{\textsc{2018}}
    &\textbf{Hemlock Research Assistant} \\
    \small{\textendash ~\textsc{2020}}
    &\footnotesize{\emph{Southern Appalachian Cooperative Ecosystems Studies Unit, U.S. Park Service}} \\
    &\footnotesize{Assessed efficacy of biological and chemical treatments of hemlock woolly adelgid by monitoring hemlock tree canopy health, hemlock woolly adelgid infestation rates, and \textit{Laricobius} spp. predator beetle densities.}
\end{longtable}